The United States stands nearly alone among established democracies in its resistance to proportional representation. While most democratic nations ensure that legislative seat shares approximate vote shares \citep{lijphart1999patterns, gallagher2011election}, the American system of single-member, winner-take-all districts routinely produces substantial disproportionality: a party receiving 55\% of votes might win 65\% of seats, or conversely, a party with 48\% support might control only 35\% of legislative seats \citep{stephanopoulos2015partisan}. This mismatch between votes and seats—termed the ``efficiency gap'' in recent scholarship \citep{stephanopoulos2015partisan}—is not merely a statistical artifact but a fundamental challenge to representative democracy.

Why does the United States maintain this system? The answer lies not in constitutional mandate but in legal and cultural path dependence. The Uniform Congressional District Act of 1967 requires that U.S. House members be elected from single-member districts \citep{ucda1967}, codifying a practice dating to 1842. This requirement forecloses the proportional representation systems used internationally: party-list PR (Netherlands, Israel), mixed-member PR (Germany, New Zealand), and single transferable vote in multi-member districts (Ireland, Malta) \citep{farrell2011electoral, massicotte2004mixed}. All these systems abandon single-member districts in favor of multi-member constituencies or compensatory seats. As a result, American electoral reformers face an apparent impasse: proportional representation requires multi-member districts, but federal law mandates single-member districts.

This paper proposes a radical solution: \emph{overlapping party-specific single-member districts}. Rather than drawing one set of districts that determines all House seats, we propose that each party independently draw districts for the seats they are allocated through proportional vote shares. If Democrats win 6 of 10 seats in a state based on statewide vote totals, they draw 6 districts covering the entire state. If Republicans win 4 seats, they draw 4 separate districts, also covering the entire state. These districts \emph{overlap}: every geographic point belongs to both a Democratic district and a Republican district. Every voter has multiple representatives—one from each party that won seats in their state.

This system is unprecedented in American politics and, to our knowledge, in any democracy worldwide. It achieves perfect proportionality by design: seat shares exactly match vote shares because seats are allocated proportionally before districting occurs. It maintains single-member districts: each party's districts elect one representative each, satisfying the 1967 Act's requirement. Yet it fundamentally transforms the relationship between geography and representation: voters no longer share districts with geographic neighbors but instead are grouped by partisan affiliation within overlapping spatial boundaries.

The implications are profound. \textbf{Representation}: Every voter gains a co-partisan representative regardless of their local geographic majority. A Democrat in rural Wyoming and a Republican in urban Manhattan—both currently in districts won by the opposite party—would each have representatives from their preferred party under this system. \textbf{Gerrymandering}: The practice becomes irrelevant, as each party controls its own districting without affecting the other party's seat count. \textbf{Third parties}: Proportional allocation gives minor parties representation if they exceed state-specific thresholds (typically 1/(N+1) where N is the state's seat count). \textbf{Accountability}: Voters face new complexity—which of their multiple representatives do they contact for constituent services? How do overlapping constituencies affect legislative behavior?

We implement this system algorithmically for all 50 states using recursive bisection \citep{author2024recursive}, a graph partitioning method that produces compact, population-balanced districts. The implementation is straightforward: allocate seats to parties using D'Hondt proportional allocation \citep{gallagher2011election}, then run independent recursive bisection for each party's allocated seats. The result is a complete set of overlapping party-specific districts for the entire nation, allowing direct comparison to both current enacted districts and neutral algorithmic districts.

Our analysis reveals both the promise and the peril of this system. \textbf{Proportionality}: Guaranteed by construction—efficiency gaps drop to exactly zero. \textbf{Compactness}: Party-specific districts remain nearly as compact as neutral algorithmic districts, suggesting geographic constraints do not substantially hinder party-based redistricting. \textbf{Representation quality}: 100\% of voters gain co-partisan representation, compared to roughly 60-70\% under current enacted districts. \textbf{Complexity}: The overlay of multiple district maps creates unprecedented administrative and cognitive burdens for voters, election officials, and candidates.

This paper makes three contributions. \textbf{Theoretical}: We expand the design space of proportional representation systems, demonstrating that single-member districts and proportional outcomes are compatible if one accepts overlapping geographies. \textbf{Algorithmic}: We provide a complete computational implementation, generating party-specific districts for all 50 states and quantifying compactness, representational coverage, and geographic patterns. \textbf{Normative}: We illuminate fundamental trade-offs in democratic representation—between geographic and partisan bases of representation, between simplicity and proportionality, between community cohesion and partisan fairness.

We proceed as follows. Section 2 details the system design: proportional seat allocation, party-specific redistricting, and overlapping geography. Section 3 describes our algorithmic implementation using recursive bisection and comparison baselines. Section 4 presents results for all 50 states, comparing proportionality, compactness, and representation quality across enacted, neutral algorithmic, and party-specific systems. Section 5 discusses democratic trade-offs, international comparisons, and implementation challenges. Section 6 concludes with implications for redistricting reform and future research directions.
